% ECONOMICS_PROBLEM: {"id": "E32-456", "title": "Fundamental sources of business cycles", "jel_code": "E32", "source_id": "OP-0011"}
% !TeX program = xelatex
\documentclass[11pt,a4paper]{article}
\usepackage[margin=23mm]{geometry}
\usepackage{fontspec}
\setmainfont{Latin Modern Roman}
\usepackage{amsmath,amssymb,mathtools}
\usepackage{xcolor,enumitem,booktabs,longtable,array,xurl,hyperref}
\definecolor{muted}{HTML}{53636C}
\hypersetup{colorlinks=true,urlcolor=blue,linkcolor=blue}
\setlength{\parindent}{0pt}
\setlength{\parskip}{5pt}
\newcommand{\R}{\mathbb R}
\newcommand{\E}{\mathbb E}
\newcommand{\N}{\mathbb N}
\newcommand{\ind}{\mathbf 1}
\newcommand{\Var}{\operatorname{Var}}
\newcommand{\Cov}{\operatorname{Cov}}
\newcommand{\supp}{\operatorname{supp}}
\DeclareMathOperator*{\argmin}{arg\,min}
\DeclareMathOperator*{\argmax}{arg\,max}
\newcommand{\entrymeta}[3]{{\sffamily\small\color{muted}\textbf{\texttt{#1}}\quad|\quad #2\par #3\par}}
\newcommand{\Problemref}[1]{\texttt{#1}}
\newcommand{\JELref}[2]{\texttt{#2} (JEL \texttt{#1})}
\begin{document}
\section*{Fundamental sources of business cycles}
\textbf{Problem ID:} \texttt{E32-456}\quad \textbf{JEL:} \texttt{E32}\par
% BEGIN ECONOMICS STATEMENT
\label{problem:OP-0011}
\label{atlas:prob:A1}

\noindent\textbf{Original atlas ID:} A1. \textbf{Formulation:} editorial specification of a research family.

\subsection*{Definitions and assumptions}

Specify an outcome vector $y_t\in\mathbb R^k$, sample period, variable definitions, and a stable linear innovation representation
\[
y_t=\sum_{h=0}^\infty\Psi_h B\varepsilon_{t-h},\qquad
 \Psi_0=I_k,\quad\mathbb E\varepsilon_t=0,\quad
 \mathbb E[\varepsilon_t\varepsilon_t^\top]=I_k,
\]
with serially uncorrelated shocks and square-summable impulse responses. The reduced-form innovation covariance is $\Sigma_u=BB^\top$. Columns $b_j$ of $B$ acquire economic labels only through declared restrictions, such as an external instrument, timing exclusion, sign restriction, or long-run restriction. For coordinate $\ell$ and forecast horizon $H\ge1$, define its forecast-error variance share, when the denominator is positive, by
\[
\omega_{\ell j}(H;B)=
\frac{\sum_{h=0}^{H-1}(e_\ell^\top\Psi_h b_j)^2}
 {\sum_{r=1}^k\sum_{h=0}^{H-1}(e_\ell^\top\Psi_h b_r)^2},
\]
where $e_\ell$ is the coordinate vector.

\subsection*{Formal research question}

Given the observable law and a declared set $\mathcal B_I$ of structural impact matrices consistent with it and the identifying restrictions, characterize the joint identified set
\[
\mathcal O_I(H)=\{(\omega_{\ell j}(H;B))_{\ell,j}:B\in\mathcal B_I\}.
\]
Which additional independently defensible measurements or instruments narrow this set enough to distinguish technology, demand, financial, markup, uncertainty, news, and policy contributions? Construct confidence regions covering the whole feasible share vector, and compare model-specific decompositions using the same outcomes, horizons, and shock normalization. If some disturbances are correlated, specify an allocation convention rather than interpreting the squared-column formula as an invariant set of additive fractions.

\subsection*{Interpretation and scope}

The shares are conditional on a structural model, sample, outcome, and horizon. With orthogonal normalized shocks they are nonnegative and sum to one for each outcome; this is a property of the declared forecast-error decomposition, not a universal causal accounting identity. Rotating $B$ can preserve $\Sigma_u$ while changing economic labels and shares. External instruments need relevance and exclusions tied to a particular shock; naming an instrument does not prove those assumptions. Large nonlinear disturbances, regime changes, and nonfundamental representations require a different decomposition. Historical contributions and unconditional business-cycle variance are different targets and should not be reported interchangeably with the displayed forecast-error shares.

\subsection*{Source audit and status}

Kydland--Prescott (1982), Galí (1999), and Smets--Wouters (2007) are verified. They respectively provide a calibrated real-business-cycle benchmark, technology-shock evidence under structural VAR restrictions, and an estimated DSGE decomposition. They already answer restricted versions of the question, with different maintained mechanisms and data. The source atlas does not identify one invariant target or restriction set, so its ``fractions'' cannot be formalized as model-free constants. The rotation-based identified set is editorial. The remaining agenda is structural discrimination and robustness across defensible models; these references alone do not certify a particular globally unresolved variance-share theorem in 2026.

\subsection*{References}

\begin{enumerate}

\item[S1] Finn E. Kydland and Edward C. Prescott (1982). \emph{Time to Build and Aggregate Fluctuations}. Econometrica 50(6), 1345--1370. \url{https://www.jstor.org/stable/1913386}. \textit{Locator:} article; calibrated stochastic growth model. \textit{Establishes:} A real-business-cycle explanation under specified technology shocks and calibration. \textit{Does not establish:} a uniquely identified empirical decomposition of all aggregate fluctuations.

\item[S2] Jordi Galí (1999). \emph{Technology, Employment, and the Business Cycle: Do Technology Shocks Explain Aggregate Fluctuations?}. American Economic Review 89(1), 249--271. \url{https://www.aeaweb.org/articles?id=10.1257/aer.89.1.249}. \textit{Locator:} abstract; structural VAR identification. \textit{Establishes:} Technology-shock evidence conditional on identifying restrictions. \textit{Does not establish:} restriction-free shock shares.

\item[S3] Frank Smets and Rafael Wouters (2007). \emph{Shocks and Frictions in US Business Cycles: A Bayesian DSGE Approach}. American Economic Review 97(3), 586--606. \url{https://www.aeaweb.org/articles?id=10.1257/aer.97.3.586}. \textit{Locator:} abstract; estimated DSGE decomposition. \textit{Establishes:} A model-specific estimated business-cycle shock decomposition. \textit{Does not establish:} a universally invariant allocation of variance across shock labels.

\end{enumerate}

\subsection*{Common conventions: Source mathematical conventions}

Unless an entry states otherwise, agent and firm sets are finite. State and action spaces are measurable, strategies and outcome maps are measurable, probability laws are specified, and expectations exist for the targets under discussion. Infinite-horizon discount factors lie in $(0,1)$ and discounted objectives are finite. Norms, tolerances and complexity input sizes must be specified for computational comparisons. Constants or primitives that appear locally are inputs, not empirical facts asserted by this catalogue.

\textbf{Identification.} A statistical model $m\in\mathcal M$ implies an observable law $P_m$ and target $\theta(m)$. At law $P$, its identified set is
\[
\Theta(P;\mathcal M)=\{\theta(m):m\in\mathcal M,\ P_m=P\}.
\]
Point identification holds when this set is a singleton, or equivalently when $P_{m_1}=P_{m_2}$ implies $\theta(m_1)=\theta(m_2)$ for all compatible models. Sharp bounds mean bounds attained, or approached when the boundary is not attained, by compatible models. Writing this set is a definition, not a solution: the substantive task is to establish informative restrictions and derive or estimate the set. An unrestricted-confounding model may give only trivial bounds.

\textbf{Inference.} Identification, estimability and valid uncertainty quantification are separate questions. Fix $\alpha\in(0,1)$. A confidence procedure $C_n$ over a declared law class $\mathcal P$ has asymptotic uniform coverage at level $1-\alpha$ when
\[
\liminf_{n\to\infty}\inf_{P\in\mathcal P}\Pr_P\{\theta(P)\in C_n\}\geq1-\alpha.
\]
For set-identified targets the coverage event must be defined separately, for example coverage of the full identified set. If an entry uses identification-indexed models instead of uniquely indexed laws, the same coverage convention is applied over those models. Pointwise limits do not establish uniform validity, and asymptotic validity does not guarantee finite-sample precision. Sample sizes, dependence structures and weak-identification sequences are part of each application.

\textbf{Counterfactuals.} $Y_i(a)$ is a potential outcome under a specified intervention, including its timing, implementation and versions. It is not $Y_i$ conditional on voluntarily choosing $a$. A full intervention vector permits interference; shorthand scalar treatment notation does not silently impose no interference. Assignment, selection, attrition, anticipation and measurement must be specified. A claim about an instrument's local effect is not automatically a population policy effect.

\textbf{Economic equilibrium.} A counterfactual model includes best responses, constraints, feasibility, market clearing, state transitions and expectations consistent with the stated information. When several equilibria exist, either a declared selector is an input or the target is set-valued. Comparisons across policy regimes must use the stated selection convention. Reduced-form relationships are not presumed invariant after an equilibrium-changing intervention.

\textbf{Optimization and welfare.} Optimization is over the stated feasible set. If attainment is not established, $\sup$ or $\inf$ and $\varepsilon$-optimal choices replace an asserted maximizer or minimizer. For example, with value $V=\sup_{a\in\mathcal A}W(a)<\infty$, an $\varepsilon$-optimal set is $\{a\in\mathcal A:W(a)\geq V-\varepsilon\}$ for $\varepsilon>0$. Benefits, costs, risk and health or environmental outcomes can be aggregated only using declared units and conversion weights. Interpersonal utility weights, the treatment of fiscal transfers and foreign profits, discounting, and the behavioral welfare criterion are explicit inputs. A comparative-static derivative additionally requires existence and appropriate differentiability of the selected counterfactual.

\textbf{Sources.} Local labels [S1], [S2], and so forth restart in every question. Locators identify the relevant abstract, model, section or result at the level actually checked; a bibliographic or abstract-level verification is not represented as inspection of an entire proof. Working papers are distinguished from later publications and changing versions. Source summaries are paraphrased. All URL checks and editorial formulations refer to the audit date stated above.
% END ECONOMICS STATEMENT
\end{document}
